% ---------------------------------------------------------------------------
% P1-statement.tex --- formal statement of Proposition \ref{prop:band}
% (interior masking band in consensus integrity). Proof: appendix/P1-proof.tex.
% Converted from docs/propositions-workdir/P1-draft.md (content frozen;
% notation renames only: kappa -> eta, theta -> r per the paper's canon).
% ---------------------------------------------------------------------------

The platform chooses a disclosure granularity $k \in [0,1]$ ($k = 1$ full disclosure,
$k = 0$ full masking). A coalition is good with probability $qc$ and bad with
probability $1 - qc$, where $q \in (0,1]$ is authentication accuracy and $c \in (0,1)$
is consensus integrity. A good coalition has robustness $r_g \sim F_c$ on $[0,1]$ and
delivers surplus $V > 0$ unless it unravels, which occurs iff $r_g < \hat r(k)$, where
$\hat r(k)$ is the strategic activation cutoff induced by disclosure level $k$; a bad
coalition imposes loss $L > 0$ unless screened, which happens with probability $s(k)$.
Welfare is
\begin{equation}
W(k; q, c) \;=\; qcV\bigl(1 - F_c(\tau k)\bigr) \;-\; (1 - qc)\,L\,(1 - k).
\tag{W}\label{eq:Wband}
\end{equation}
Let $K^*(c) := \operatorname*{arg\,max}_{k \in [0,1]} W(k; q, c)$ (the dependence on
$q$ is suppressed except in part~(e)). We write ``full disclosure is uniquely optimal
at $c$'' for $K^*(c) = \{1\}$ and ``masking is optimal at $c$'' for $1 \notin K^*(c)$.
Define the \emph{masking set}
\[
M \;:=\; \{\, c \in (0,1) : 1 \notin K^*(c) \,\},
\]
and the \emph{marginal-balance function} (minus the right-hand slope of $W$ at $k = 1$)
\begin{equation}
\varphi(c) \;:=\; qcV\tau f_c(\tau) \;-\; (1 - qc)\,L.
\tag{$\varphi$}\label{eq:phi}
\end{equation}
$\varphi(c) > 0$ says: at full disclosure, the marginal unraveling loss
$qcV\tau f_c(\tau)$ exceeds the marginal screening gain $(1 - qc)L$.

\begin{assumption}\label{ass:band}
\leavevmode
\begin{itemize}
\item[\textbf{(A1)}] \textbf{(Payoffs.)} $V, L > 0$, $\tau \in (0,1)$, $q \in (0,1]$,
  $c \in (0,1)$; welfare is \eqref{eq:Wband}.
\item[\textbf{(A2)}] \textbf{(Linear channels --- restriction.)} $\hat r(k) = \tau k$,
  $s(k) = k$, hence the unraveling risk is $\rho(k,c) = F_c(\tau k)$. This is exactly
  the specification of the numerical ground-truth scripts \texttt{v3.py}~/
  \texttt{v3\_shape.py}.
\item[\textbf{(A3)}] \textbf{(Family regularity.)} For each $c \in (0,1)$, $F_c$ is a
  CDF on $[0,1]$ with density $f_c$ strictly positive and continuous on $(0,1)$, and
  $(r, c) \mapsto f_c(r)$ is continuous on $(0,1) \times (0,1)$.
\item[\textbf{(A4)}] \textbf{(Fragile-end thinness at the cutoff.)} There exists
  $\bar k \in (0,1)$ with
  \[
  \lim_{c \to 0^+} \; c \, \sup_{r \in [\tau\bar k,\, \tau]} f_c(r) \;=\; 0.
  \]
  (At very low integrity, the robustness density of good coalitions carries vanishing
  $c$-weighted mass near the transparency cutoff $\tau$.)
\item[\textbf{(A5)}] \textbf{(Robust-end thinness below the cutoff.)}
  {\sloppy
  $\lim_{c \to 1^-} S(c) = 0$, where $S(c) := \sup_{r \in (0,\tau]} f_c(r)$, and
  \textbf{either} (i) $q < 1$, \textbf{or} (ii)
  \[
  \limsup_{c \to 1^-} \; \frac{cV\tau S(c)}{(1 - c)L} \;<\; 1.
  \]
  (At very high integrity, good coalitions are too robust to unravel: the density on
  $[0, \tau]$ vanishes --- fast enough, relative to $1 - c$, if $q = 1$.)\par}
\item[\textbf{(A6)}] \textbf{(Interior bite --- non-vacuity.)} There exists
  $c_m \in (0,1)$ with $\varphi(c_m) > 0$, i.e.
  \[
  q\, c_m V \tau f_{c_m}(\tau) \;>\; (1 - q c_m)\, L.
  \]
\item[\textbf{(A7)}] \textbf{(Beta specialization, for parts (d).)}
  $F_c = \mathrm{Beta}\bigl(c\eta, (1-c)\eta\bigr)$ on $[0,1]$ with fixed concentration
  $\eta > 2$ (mean $c$; the family of v2/v3 and of \texttt{v3\_shape.py}). Define
  \[
  \bar c \;:=\; \frac{1 + \tau(\eta - 2)}{\eta} \;\in\; \Bigl(\frac{1}{\eta},\, 1\Bigr)
  \qquad \text{(the concavity threshold)},
  \]
  the smallest $c$ at which the Beta mode $(c\eta - 1)/(\eta - 2)$ reaches $\tau$.
\end{itemize}
\end{assumption}

\begin{proposition}[Interior masking band in $c$]\label{prop:band}
Assume A1--A3.
\begin{itemize}
\item[\textup{(a)}] \textup{(Existence.)} For every $c \in (0,1)$, $W(\cdot\,; c)$ is
continuous on $[0,1]$ and $K^*(c) \neq \emptyset$.
\end{itemize}
Assume in addition A4--A6. Then:
\begin{itemize}
\item[\textup{(b)}] \textup{(Full disclosure at both extremes.)} There exist
$0 < c_1 \le c_2 < 1$ such that $K^*(c) = \{1\}$ for every $c \in (0, c_1)$ and every
$c \in (c_2, 1)$.
\item[\textup{(c)}] \textup{(Masking in the middle; the band.)} $1 \notin K^*(c_m)$, so
$M \neq \emptyset$. Consequently
\[
c_{\mathrm{lo}} := \inf M \quad\text{and}\quad c_{\mathrm{hi}} := \sup M
\]
satisfy $0 < c_1 \le c_{\mathrm{lo}} \le c_m \le c_{\mathrm{hi}} \le c_2 < 1$; full
disclosure is optimal ($1 \in K^*(c)$) for every
$c \in (0, c_{\mathrm{lo}}) \cup (c_{\mathrm{hi}}, 1)$, uniquely so on
$(0, c_1) \cup (c_2, 1)$; and masking is strictly optimal at
$c_m \in [c_{\mathrm{lo}}, c_{\mathrm{hi}}]$. In particular the masking region is
contained in a band strictly interior to $(0,1)$ and is nonempty.
\end{itemize}
Assume in addition A7. Then:
\begin{itemize}
\item[\textup{(d)}] \textup{(Sharp upper threshold and FOC characterization on the
upper shoulder.)}
\begin{itemize}
\item[\textup{(d1)}] $c \mapsto \log\bigl[\, qcV\tau f_c(\tau) \,/\, ((1-qc)L) \,\bigr]$
is strictly concave on $(0,1)$; hence $\{c : \varphi(c) > 0\}$ is an open interval
$(c_-, c_+)$, nonempty by A6, with $c_- < c_m < c_+$.
\item[\textup{(d2)}] For every $c \in [\bar c, 1)$, $W(\cdot\,; c)$ is strictly concave
on $[0,1]$, $K^*(c)$ is a singleton $\{k^*(c)\}$, $k^*(c) > 0$, and
\[
k^*(c) < 1 \quad\iff\quad \varphi(c) > 0;
\]
when $\varphi(c) > 0$, $k^*(c) \in (0,1)$ is the unique solution of the first-order
condition
\[
qcV\tau f_c(\tau k^*) = (1 - qc)L
\qquad \text{(marginal unraveling loss $=$ marginal screening gain).}
\]
\item[\textup{(d3)}] If moreover $c_+ > \bar c$ (true at the script parameters
$(V, L, \tau, \eta) = (1, 1.2, 0.7, 8)$ with $q = 0.95$; Remark~\ref{rem:anchors}),
then $c_{\mathrm{hi}} = c_+ < 1$, the interval
$\bigl( \max(\bar c, c_-),\, c_+ \bigr) \subseteq M$, $k^*$ is continuous on
$\bigl( \max(\bar c, c_-),\, c_+ \bigr)$, and $k^*(c) \to 1$ as
$c \uparrow c_{\mathrm{hi}}$: the upper transition is a continuous closing of the band,
characterized exactly by $\varphi(c_{\mathrm{hi}}) = 0$.
\end{itemize}
\item[\textup{(e)}] \textup{(Comparative statics in $q$; needs only A1--A3.)} For
$q < q' \le 1$:
\begin{itemize}
\item $M(q) \subseteq M(q')$: the masking set expands with authentication accuracy;
hence, whenever $M(q) \neq \emptyset$ (so that the thresholds $c_{\mathrm{lo}}(q),
c_{\mathrm{hi}}(q)$ are defined --- and then $M(q') \neq \emptyset$ too, by the
inclusion), $c_{\mathrm{lo}}(\cdot)$ is non-increasing and $c_{\mathrm{hi}}(\cdot)$
non-decreasing in $q$. (When $M(q) = \emptyset$ the thresholds are undefined ---
$\inf \emptyset$~/ $\sup \emptyset$ --- and the clause is vacuous.)
\item $\max K^*(c; q') \le \min K^*(c; q)$ for every $c$: every optimal disclosure
level under $q'$ is weakly below every optimal level under $q$. In particular, wherever
$K^*$ is a singleton, $k^*(c; q)$ is non-increasing in $q$ (more masking as
authenticity rises).
\end{itemize}
\end{itemize}
\end{proposition}

\begin{conjecture}[Connectedness of $M$]\label{conj:connected}
$M$ is an interval, i.e.\ $M = (c_{\mathrm{lo}}, c_{\mathrm{hi}})$ exactly.
\end{conjecture}
